% results.tex (국문 v14, 2026-10-20): 영문 v14 와 같은 구성. 간결한 캡션(세부는 방법 절), Fig. 6 문단은 v14 설계. Δ = RMSE(앞) − RMSE(뒤), 같은 지점-연도; 괄호 안은 95% 짝지은 블록 부트스트랩 구간.
\section{결과}\label{sec:results}

\subsection{관측 연도 기후 입력의 효과}
관측은 시간적으로 고르지 않다(그림~\ref{fig:data}b, 표~\ref{tab:data}). 알래스카 지점-연도의 \ntopyrshareAk\%가 \ntopyrAk{}년에 관측된 반면, 레나 델타 라벨은 \nnyrLe{}개 해에 퍼져 있다. 측정은 대부분 융해 계절이 끝나기 전인 8월에 이루어졌고(관측일 중앙값: 알래스카 \ndoymedAk{}일, 레나 델타 \ndoymedLe{}일, 캐나다 \ndoymedCa{}일), 관측일까지 누적한 TDD는 계절 전체 합의 중앙값 \ntddfracAk, \ntddfracLe, \ntddfracCa{}에 해당하였다(식~(\ref{eq:tddobs})).

평년 기후 입력을 관측 연도 입력으로 바꾸면 재보정 Stefan 모델의 공간 블록 RMSE가 레나 델타에서($\Delta$RMSE \ndforcLePo{} cm), 물리 잔차 모델의 RMSE가 레나 델타(\ndforcLeRo{} cm)와 캐나다(\ndforcCaRo{} cm)에서 줄었다(그림~\ref{fig:accuracy}a). 알래스카에서는 변화가 판정되지 않았다(Stefan \ndforcAkPo{} cm). 라벨이 한 해에 몰려 있어 관측 연도 입력이 공간 분포 이상의 정보를 거의 더하지 않기 때문이다. 정확도와 별개로 관측 연도 입력은 각 예측을 명시된 해의 기후에 묶으며, 아래의 연도별 지도는 이를 전제로 한다.

\paperfig{Fig3_accuracy}{1.0}{\textbf{기후 입력, 검증 설계, 물리 잔차 이득의 조건.} \textbf{a}, 2011--2024년 평균 기후 입력을 관측 연도 입력으로 바꿀 때의 공간 블록 RMSE 변화. \textbf{b}, 다섯 가지 검증 설계에서 네 방법의 RMSE(위, 축 밖 값은 숫자로 표기)와 각 설계에서 물리 잔차 모델과 재보정 Stefan 모델의 차이(아래). \textbf{c}, 공간 블록 검증에서 공변량 묶음 여섯 가지에 대한 같은 차이. \textbf{d}, 공변량 여덟 개의 지역 안 3분위에서의 같은 차이. 테두리 칸은 구간이 0을 제외하고, 빗금 칸은 지점-연도 30개 또는 블록 5개 미만이다. 구간은 95\% 짝지은 블록 부트스트랩 구간이며, 음수는 앞에 적은 모델이 낫다는 뜻이다.}{fig:accuracy}

\subsection{공간 · 시간 · 지역 검증에서의 정확도}
모델 순위는 검증 설계에 따라 달랐다(그림~\ref{fig:accuracy}b, 보충 표 S5). 무작위 교차검증은 직접 ML에 가장 낮은 오차를 주었고(알래스카 \nrmseAkrandDz{} cm), 공간 블록 교차검증은 직접 ML의 오차를 알래스카 \noptimAkDz, 레나 델타 \noptimLeDz, 캐나다 \noptimCaDz{} cm 키웠다. 연도 홀드아웃은 레나 델타의 직접 ML 오차를 \nrmseLeblkDz{} cm에서 \nrmseLeyrDz{} cm로 키웠고, 블록 × 연도 설계는 지역 라벨로 학습한 모든 방법에서 지역 안 최대 오차를 냈다. 한 지역 전체를 빼면 레나 델타의 직접 ML 오차는 \nrmseLeregDz{} cm까지 올랐고, 다른 지역 계수를 쓴 Stefan 모델은 \nrmseLeregPz{} cm였다($\Delta$RMSE \ndrefLeregDz{} cm). 같은 시험에서 물리 잔차 모델은 Stefan 모델 수준을 유지하였고(\ndrefLeregRo{} cm) 캐나다에서는 Stefan 모델을 개선하였다(\ndrefCaregRo{} cm). 따라서 공간 블록 교차검증을 주 설계로, 지역 홀드아웃을 전이성 시험으로 쓴다.

\subsection{물리 잔차 학습의 이득과 그 조건}
관측 연도 입력의 공간 블록 교차검증에서 물리 잔차 모델은 재보정 Stefan 모델의 RMSE를 레나 델타에서 \nrmseLeblkPo{} cm에서 \nrmseLeblkRo{} cm로($\Delta$RMSE \ndrefLeblkRo{} cm), 캐나다에서 \nrmseCablkPo{} cm에서 \nrmseCablkRo{} cm로(\ndrefCablkRo{} cm) 줄였으며, 알래스카에서는 두 모델이 \nrmseAkblkPo{} cm와 \nrmseAkblkRo{} cm였다(\ndrefAkblkRo{} cm). 이득은 잔차 학습기가 쓸 수 있는 공변량에 좌우되었다(그림~\ref{fig:accuracy}c). 기후 공변량만으로는 잔차 모델이 알래스카(\ndfeatAkFclim{} cm)와 캐나다(\ndfeatCaFclim{} cm)에서 Stefan 모델을 개선하지 못하였고, 지형 · 토양 공변량을 더하자 위의 이득이 나타났다. MODIS 식생 · 눈 시기나 고해상 토양 · 피복 분류는 레나 델타와 캐나다에서 결과를 0.5 cm 미만으로 바꾸었고, 후자는 알래스카에서 작은 이득을 더하였다(\ndfeatAkFhires{} cm).

결과는 물리 모델이 학습기에 들어가는 방식에도 좌우되었다(보충 그림 S2). Stefan 유사라벨은 Stefan 모델을 재현하였고(레나 델타 \ndrefLeblkDo{} cm), Stefan 예측을 입력 특징으로 쓰면 직접 ML과 비슷한 결과가 나왔다. 한 지역의 라벨을 모두 쓰면 직접 ML의 점 오차가 물리 잔차 모델보다 낮았으나(알래스카 \nrmseAkblkDz{} 대 \nrmseAkblkRo, 레나 델타 \nrmseLeblkDz{} 대 \nrmseLeblkRo, 캐나다 \nrmseCablkDz{} 대 \nrmseCablkRo{} cm), 재보정 Stefan 모델과의 차이는 캐나다에서만 판정되었고(\ndrefCablkDz{} cm) 지역 홀드아웃에서는 오차가 커졌다. CatBoost 를 랜덤 포레스트, 다층 퍼셉트론, TabM, FT-Transformer 로 바꾸어도 레나 델타와 캐나다에서 물리 잔차 이득의 부호는 바뀌지 않았다.

공변량 여덟 개의 지역 안 3분위에서 이득은 레나 델타(구간이 0을 제외한 3분위 \ngcLesig{}/\ngcLecells{})와 캐나다(\ngcCasig{}/\ngcCacells)의 대부분 조건에서 나타났고 Stefan 모델의 오차와 함께 커졌다(3분위 간 상관: 레나 델타 \ngcLecorr, 캐나다 \ngcCacorr; 그림~\ref{fig:accuracy}d). 레나 델타의 가장 따뜻한 TDD 3분위에서 이득은 \ndtddLewarm{} cm에 이르렀고, 알래스카에서는 어느 3분위도 \ngcAkmaxabs{} cm 넘게 바뀌지 않았다. 잔차 모델의 평균 절대 SHAP 값에서 지형 변수의 몫은 레나 델타에서 \nshapLeResTerrain\%, 토양 변수의 몫은 캐나다에서 \nshapCaResSoil\%였다(보충 그림 S3).

\subsection{표본 효율과 새 지역으로의 전이}
한 지역의 지점-연도에서 뽑은 학습 라벨 25개로 물리 잔차 모델은 같은 라벨로 재보정한 Stefan 모델 이상이었고(알래스카 \nllAkRotwfive, 레나 델타 \nllLeRotwfive, 캐나다 \nllCaRotwfive{} cm), 100개에서는 세 지역 모두에서 이득이 판정되었다(\nllAkRohundred, \nllLeRohundred, \nllCaRohundred{} cm; 그림~\ref{fig:labels}a). 라벨 25개로 학습한 직접 ML은 알래스카(\nllAkDztwfive{} cm)와 캐나다(\nllCaDztwfive{} cm)에서 Stefan 모델보다 점 오차가 컸다. 물리 입력 ML과 물리 유사라벨은 전체 자료 결과의 양상을 따랐다(보충 그림 S4).

\paperfig{Fig4_labels}{1.0}{\textbf{적은 라벨, 새 지역, 외부 CALM 지역에서의 정확도.} \textbf{a}, 각 지역에서 같은 $n$개 학습 라벨로 재보정한 Stefan 모델 대비 각 방법의 RMSE 차이(무작위 3회 평균; 띠는 물리 잔차 모델의 95\% 블록 부트스트랩 구간). \textbf{b}, 다른 두 지역과 대상 지역 라벨 $n$개로 학습한 모델의 지역별 RMSE. 축 밖 값은 숫자로 표기하였다. \textbf{c}, 소지역 번호를 붙인 외부 CALM 지점과 연구 지역의 라벨 셀. \textbf{d}, 세 제품 값이 모두 있는 외부 지점-연도 \nftrows{}개에서 같은 지역의 라벨 수에 따른 RMSE. 수평선은 제품이다. \textbf{e}, 소지역별로 지역 라벨 0개와 5개일 때 물리 잔차 모델의 RMSE와 ESA CCI v5의 RMSE(괄호 안은 지점-연도 수).}{fig:labels}

전이 실험에서는 각 지역을 다른 두 지역과 대상 지역 라벨 $n$개로 학습한 모델로 예측하였다(그림~\ref{fig:labels}b). 대상 라벨 없이 직접 ML은 레나 델타에서 \ntlrmseLeDszero{} cm였고, 원천 계수 Stefan 모델은 \ntlrmseLePzzero{} cm, 물리 잔차 모델은 \ntlrmseLeRszero{} cm였다. 대상 라벨 25개는 물리 잔차 오차를 \ntlrmseLeRstwfive{} cm로, 직접 ML 오차를 \ntlrmseLeDstwfive{} cm로 줄였다. 알래스카에서는 물리 잔차 모델이 모든 $n$에서 가장 낮은 오차를 냈다(대상 라벨 없이 \ntlrmseAkRszero{} cm). 캐나다에서는 원천 지역으로 학습한 직접 ML이 가장 정확하였다(\ntlrmseCaDszero{} cm). 캐나다 라벨 25개로 Stefan 계수를 재보정하면 오차가 \ntlrmseCaPzzero{} cm에서 \ntlrmseCaPntwfive{} cm로 커졌고, 이 기준값을 쓰는 물리 잔차 모델도 뒤따랐다(\ntlrmseCaRszero{} cm에서 \ntlrmseCaRstwfive{} cm). 캐나다 지점의 융해 깊이 범위가 넓기 때문이다.

같은 절차를 모델 개발에 쓰지 않은 지점-연도 \nftScoredRows{}개의 외부 CALM 8개 지역에 적용하였다(그림~\ref{fig:labels}c--e, 보충 그림 S5). Wei 등, Liu 등, ESA CCI v5 값이 있는 7개 지역의 지점-연도 \nftrows{}개에서, 세 연구 지역으로 학습한 물리 잔차 모델은 지역 라벨 없이 \nftAllRszero{} cm였다. 같은 지역의 라벨 5개는 오차를 \nftAllRsfive{} cm로, 25개는 \nftAllRstwfive{} cm로 줄였고, 직접 ML은 라벨 25개로 \nftAllDstwfive{} cm에 이르렀다. 라벨 5개의 물리 잔차 모델은 ALT 관측으로 학습하지 않은 ESA CCI v5보다 정확하였다(같은 지점-연도에서 \nftAllCcifivezero{} cm; ESA CCI v5 값이 있는 지점-연도 \nftdCciRows{}개 전체에서 $\Delta$RMSE \nftdRsCcifivefive{} cm). 이 지점들을 학습에 쓴 Wei 등과 Liu 등은 \nftAllWeizero{} cm와 \nftAllLiuzero{} cm였다. 감소의 대부분은 Stefan 계수 재보정에서 왔고, 잔차 학습기는 이 지점들에서 오차를 더 줄이지 못하였다(라벨 5개에서 \nftdRsPnfive{} cm). 서시베리아 · 중앙시베리아 · 스발바르의 활동층이 연구 지역의 범위를 넘기 때문이다. 합동 곡선은 서로 다른 지역 반응을 합친 것이다(그림~\ref{fig:labels}e). 라벨 5개로 서시베리아, 중앙시베리아, 유럽 러시아, 스발바르에서는 오차가 줄었지만, 지역 라벨 없는 오차가 이미 작았던 추코트카(\nftChRszero{}에서 \nftChRsfive{} cm), 캄차카(\nftKaRszero{}에서 \nftKaRsfive{} cm; 지점-연도 16개, 미표시), 북동시베리아(\nftNsRszero{}에서 \nftNsRsfive{} cm)에서는 커졌다.

\subsection{같은 셀 · 같은 해의 공개 제품 비교}
공통 지점-연도(\nyrscc{}년)에서 물리 잔차 모델은 레나 델타에서 \nccLeRoall{} cm였고, ESA CCI v5 \nccLeCcifiveall{} cm, Aalto 등 \nccLeAaltoall{} cm, ESA CCI v4 \nccLeCcifourall{} cm, Liu 등 \nccLeLiuall{} cm, Wei 등 \nccLeWeiall{} cm였다(그림~\ref{fig:products}a, 보충 표 S6). Wei 등의 학습 지점에서 25 km보다 먼 레나 델타 지점-연도(Wei 등 값이 있는 \nnprodLeweimask{}개, ESA CCI v5 값이 있는 \nnprodLeccifivemask{}개)에서 짝지은 차이는 Wei 등 대비 \ndprodLeweiRomask{} cm, ESA CCI v5 대비 \ndprodLeccifiveRomask{} cm였다. 알래스카(\nccAkRoall{} 대 Wei 등 \nccAkWeiall, Liu 등 \nccAkLiuall{} cm)와 캐나다(\nccCaRoall{} 대 \nccCaWeiall, \nccCaLiuall{} cm)에서는 모델과 두 ML 제품이 같은 수준이었다.

레나 델타 지점-연도에서 두 ML 제품은 관측 ALT와 무관하게 대부분의 셀에 약 40--70 cm를 부여한다(그림~\ref{fig:products}b). 관측과 예측의 상관은 Wei 등 \nrLeWei, Liu 등 \nrLeLiu{}였고 물리 잔차 모델은 \nrLeRo{}였으며, 평균 편차는 \nbiasLeAllWei{} cm와 \nbiasLeAllLiu{} cm 대 \nbiasLeAllRo{} cm였다. 모든 계열의 편차는 얇은 활동층 쪽으로 커졌다(그림~\ref{fig:products}c). 관측 ALT 30 cm 미만인 지점-연도(레나 델타 \nthinNLe{}개)에서 Wei 등과 Liu 등은 \nbiasLeThinWei{} cm와 \nbiasLeThinLiu{} cm, 물리 잔차 모델은 \nbiasLeThinRo{} cm 과대 추정하였고, 알래스카의 같은 구간(\nthinNAk{}개)에서는 각각 \nbiasAkThinWei, \nbiasAkThinLiu, \nbiasAkThinRo{} cm 과대 추정하였다.

\paperfig{Fig5_products}{1.0}{\textbf{같은 셀 · 같은 해의 공개 제품 비교.} \textbf{a}, 공통 지점-연도(\nyrscc{}년)에서 세 방법과 제품 5종의 RMSE. 채운 기호는 모든 셀, 빈 기호는 Wei 등의 학습 지점에서 25 km보다 먼 셀이며, 45 cm를 넘는 값은 압축한 구간에 그렸다. \textbf{b}, 레나 델타 공통 지점-연도 \nccLeN{}개의 관측 대 예측 ALT(1:1 선, Pearson 상관). \textbf{c}, 관측 ALT 구간별 평균 편차(축 아래는 구간별 지점-연도 수, 10개 미만인 구간은 생략). \textbf{d}, 지점-연도 30개 이상인 해의 연도별 RMSE와 연도별 지점-연도 수. 우리 방법은 모두 공간 블록 겹 밖 예측이다.}{fig:products}

지역별 대비는 제품의 학습 자료를 따른다. 알래스카 라벨 셀은 Wei 등의 학습 지점에서 중앙값 \ndistweiAkmed{} km 떨어져 있고 Liu 등은 같은 알래스카 · 캐나다 현장 자료의 이전 판으로 학습한 반면, 레나 델타 셀은 중앙값 \ndistweiLemed{} km 떨어져 있다(보충 그림 S6). 두 제품 모두 델타 북부의 얇은 활동층을 과대 추정하며, 북위 72$^\circ$ 이북의 평균 편차는 \nbiasLeNorthWei{} cm와 \nbiasLeNorthLiu{} cm이다. 레나 델타 라벨 없이 학습한 모델도 같은 지점-연도에서 이미 Wei 등보다 정확하였으므로(\ntlvprmseLeRsWeizero{} cm 대 \ntlvprmseLeWeizero{} cm, $\Delta$RMSE \ntlvpLeRsWeizero{} cm), 이 지역의 우위는 지역 라벨에만 의존하지 않는다.

지역별로 관측 셀이 가장 많은 해(알래스카 · 캐나다 2013년, 레나 델타 2016년)의 지도는 이 차이를 공간적으로 보여 준다(그림~\ref{fig:prodmaps}). 알래스카 관측 셀 \nfeAkN{}개에서 세 지도의 RMSE는 비슷하였다(\nfeAkOof, \nfeAkWei, \nfeAkLiu{} cm). 관측과의 상관은 우리 모델(\nfeAkROof)이 제품(\nfeAkRWei, \nfeAkRLiu)보다 낮았으나, 이 중 \nfeAkNNear{}개 셀은 Wei 등의 학습 지점에서 25 km 안에 있고, 그 밖의 \nfeAkNFar{}개 셀에서는 세 지도의 오차가 비슷하였다(\nfeAkFarOof, \nfeAkFarWei, \nfeAkFarLiu{} cm; $r$ \nfeAkFarROof, \nfeAkFarRWei, \nfeAkFarRLiu). 지도 간 차이는 고도 600 m 이상의 산지에서 가장 컸고, 북위 68--70$^\circ$의 산지에서 제품은 우리 지도보다 \nmapAkFootMtnWei{} cm(Wei 등), \nmapAkFootMtnLiu{} cm(Liu 등) 깊었다. 레나 델타(2016년 관측 셀 \nfeLeN{}개)에서는 두 제품이 관측보다 깊었다(평균 편차 \nfeLeBiasWei{} cm, \nfeLeBiasLiu{} cm; RMSE \nfeLeWei{} cm, \nfeLeLiu{} cm 대 \nfeLeOof{} cm). 캐나다(북서부)에서는 관측 셀 \nfeCaN{}개의 RMSE가 우리 모델 \nfeCaOof{} cm, Wei 등 \nfeCaWei{} cm(평균 편차 \nfeCaBiasWei{} cm), Liu 등 \nfeCaLiu{} cm였고, Liu 등의 학습 자료에는 이 관측의 이전 판이 포함된다. Wei 등의 학습 지점에서 25 km 밖의 \nfeCaNFar{}개 셀에서는 \nfeCaFarOof, \nfeCaFarWei, \nfeCaFarLiu{} cm였다.

\paperfig{Fig6_product_maps}{1.0}{\textbf{지역별로 관측 셀이 가장 많은 해의 물리 잔차 모델 지도와 공개 제품 지도, 그리고 관측값.} \textbf{a--c}, 알래스카 2013년의 물리 잔차 모델, Wei 등, Liu 등 지도. 제품 지도는 우리 지도의 셀에만 표시하고, 관측 셀은 관측 ALT를 같은 색 눈금으로 채운 원으로 그렸다. \textbf{d}, 세 지도 모두 값이 있는 관측 셀 \nfeAkN{}개의 관측 대 예측 ALT(1:1 선). 채운 표지는 Wei 등의 학습 지점에서 25 km 안의 셀, 빈 표지는 그 밖의 셀이며, RMSE는 전체 셀과 빈 표지 셀에 대해 적었다. 우리 예측은 공간 블록 겹 밖 값이다. \textbf{e--h}, 레나 델타 2016년의 같은 비교(\nfeLeN{}개 셀). \textbf{i--l}, 캐나다(북서부) 2013년의 같은 비교(\nfeCaN{}개 셀). 지역마다 적용 가능 영역 안 세 지도 값의 5--95 백분위수를 색 범위로 쓰며, 옅은 색은 영역 밖의 셀이다.}{fig:prodmaps}

연도별로 물리 잔차 모델은 지점-연도 30개 이상인 레나 델타 \nywLen{}개 해(\nyrsywLeko) 모두에서 Wei 등과 Stefan 모델보다 오차가 낮았고, 알래스카에서는 Wei 등 대비 \nywAkRoWein{}개 해 중 \nywAkRoWeiwin{}개, Stefan 모델 대비 \nywAkRoPon{}개 해 중 \nywAkRoPowin{}개(\nyrsywAkko), 캐나다에서는 Stefan 모델 대비 \nywCaRoPon{}개 해 중 \nywCaRoPowin{}개에서 낮았다(그림~\ref{fig:products}d). 여러 해 관측된 알래스카 셀에서 관측과 예측의 연도 간 변화 상관은 물리 잔차 모델 \nrdeltaAkRo, Wei 등 \nrdeltaAkWei, Liu 등 \nrdeltaAkLiu{}였고, 레나 델타에서는 제품이 연도 간 변화를 더 가깝게 따랐다(\nrdeltaLeWei, \nrdeltaLeLiu{} 대 \nrdeltaLeRo).

잔차 학습기는 제품을 기준값으로 쓸 수도 있다(식~(\ref{eq:prod})). 레나 델타 라벨로 Wei 등의 재보정 제품에 잔차 학습기를 더하면 오차가 \npadLeRpPcal{} cm 바뀌었고, Stefan 기준 모델이 제품 기준 모델보다 여전히 정확하였다(\npadLeRoRp{} cm; 보충 그림 S7). 알래스카와 캐나다에서는 두 기준값을 구별할 수 없었다.

\subsection{연도가 명시된 1 km ALT 지도, 2011--2024년}
각 지역의 모든 지점-연도로 학습한 물리 잔차 모델을 2011--2024년 각 해의 ERA5-Land 기후로 적용하였다(그림~\ref{fig:maps}). 지도는 알래스카 \nncellmapAk{}개, 레나 델타 \nncellmapLe{}개, 캐나다(북서부) \nncellmapCa{}개 셀이다. 2011--2024년 평균 대비 지역 평균 편차는 알래스카에서 \nanminyrAk{}년 \nanminAk{} cm부터 \nanmaxyrAk{}년 \nanmaxAk{} cm까지, 레나 델타에서 \nanminyrLe{}년 \nanminLe{} cm부터 \nanmaxyrLe{}년 \nanmaxLe{} cm까지, 캐나다에서 \nanminyrCa{}년 \nanminCa{} cm부터 \nanmaxyrCa{}년 \nanmaxCa{} cm까지였으며, 공간 분포도 해마다 달랐다(그림~\ref{fig:maps}f). 레나 델타에서 관측 셀의 평균 편차는 지점-연도 50개 이상인 \nanobsnLe{}개 해(\nanobsyrsLe{}년) 가운데 \nanobssameLe{}개 해에서 예측 지역 편차와 부호가 같았고($r$ = \nanobsrLe), 관측 편차가 \nanobsLeTwentyObs{} cm(지점-연도 \nanobsLeTwentyN{}개)에 이른 2020년을 빼도 상관은 \nanobsrLeNoTwenty{}였다. 해마다 관측 셀 구성이 바뀌는 알래스카에서는 일치가 약하였고($r$ = \nanobsrAk) 예측의 연도 간 범위(\nanpredrangeAk{} cm)는 관측 셀 평균의 범위(\nanobsrangeAk{} cm)보다 작았다. 캐나다에서 \nanobsnCa{}개 해의 상관은 \nanobsrCa{}였으며, 2023년과 2024년(각 지점-연도 80개)의 관측 편차는 예측 범위를 크게 넘었다(그림~\ref{fig:maps}d). 연도별 지도의 선형 추세는 알래스카, 레나 델타, 캐나다 셀 평균 10년당 각각 \ntrendAk{}, \ntrendLe{}, \ntrendCa{} cm였고, 양의 추세인 셀은 각각 \ntrendposAk\%, \ntrendposLe\%, \ntrendposCa\%였다(그림~\ref{fig:maps}e). 이 추세는 모델을 거쳐 전달된 그 기간의 ERA5-Land 기후를 반영하며 관측 추세가 아니다. 적용 가능 영역은 알래스카 셀의 \naoaAk\%, 레나 델타 셀의 \naoaLe\%, 캐나다 셀의 \naoaCa\%를 포함하며, 영역 밖의 셀은 주로 산지와 영구동토 지대의 남쪽 가장자리에 있다(보충 그림 S8). 90\% 예측 구간의 반폭은 \nhwAk{}, \nhwLe{}, \nhwCa{} cm이다. 잔차 학습기가 더하는 보정량은 대부분 $\pm$10 cm 안에 있고(보충 그림 S9), 장기 관측 지점의 시계열은 보충 그림 S10에 나타내었다.

\paperfig{Fig7_maps}{1.0}{\textbf{연도가 명시된 1 km ALT 지도, 2011--2024년.} \textbf{a--c}, 알래스카, 레나 델타, 캐나다(북서부) 연도별 지도의 2011--2024년 평균 ALT. 오른쪽에 창을 1 km 셀 그대로 확대하였고, 관측 셀은 관측 평균 ALT로 채웠으며, 옅은 색은 적용 가능 영역 밖의 셀이다. \textbf{d}, 지도의 지역 평균 편차와 지점-연도 50개 이상인 해의 관측 셀 평균 편차(같은 해의 예측 편차에 맞추어 중심을 옮김; 지역마다 y 범위가 다름). \textbf{e}, 연도별 지도의 선형 추세(셀별 최소제곱). 해마다의 기후에 대한 모델 응답을 기술한다. \textbf{f}, 2011--2024년 평균 대비 알래스카의 연도별 ALT 편차.}{fig:maps}
